% v10 appendices (docs/mist_v10_plan.zh-CN.md §4.10). All numbers are \num{v10:...} macros from scripts/make_tables_v10.py.

\section{v10 Training and Checkpoint Selection}
\label{app:v10train}

\paragraph{Stage 1 (strict baseline).} The deadline rows of the two v9 directions are interleaved into
\num{v10:c:j:rows} rows such that every global batch of 64 holds 32 Zh$\rightarrow$En and 32 En$\rightarrow$Zh rows, each
in the sample order of its v9 run; at joint step $2k$ each direction has seen exactly the rows of its v9 run at step $k$,
and the learning rate (cosine from $10^{-4}$ to $10^{-5}$ over \num{v10:c:j:s1steps} steps, 100 warm-up steps) matches the
v9 schedule at the same point. LoRA, projector and new rows as in Table~\ref{tab:hyper}; 8 H200 GPUs, micro-batch 8, no
gradient accumulation. There is no checkpoint selection: the strict baseline is the last step.

\paragraph{Stage 2 (latent write alignment).} From the merged stage-1 model, Eq.~\eqref{eq:mstep} with $\tau=0.5$, new
LoRA adapters ($r=32$, $\alpha=64$), the same posterior perturbation and the same learning-rate schedule as stage 1; each
step draws about 32 read rows per direction (64 in total). Checkpoints are saved every 500 steps. From step 3{,}000 on,
each checkpoint's text is decoded on dev5 in both directions. The rule, written before stage 2 started: a checkpoint is
feasible if its dev5 text LAAL is within \num{v10:c:sel:tol}\,s of the reference system's in both directions (the
separately trained single-direction latent-write systems, $L^{\mathrm{ref}}=$ \num{v10:s2:refL:z} /
\num{v10:s2:refL:e}\,s); $J$ is the mean of the two directions' text tcpBLEU; training stops after two evaluated
checkpoints without an improvement of $J$ by more than 0.2, and the feasible checkpoint with the highest $J$ is selected.
The held-out-loss tie-break could not be computed next to training (memory) and was not needed.

\begin{table}[t]
\centering\small
\setlength{\tabcolsep}{3.5pt}
\resizebox{\columnwidth}{!}{\begin{tabular}{rrrrrrc}
\toprule
 & \multicolumn{2}{c}{Zh$\rightarrow$En} & \multicolumn{2}{c}{En$\rightarrow$Zh} & & \\
\cmidrule(lr){2-3}\cmidrule(lr){4-5}
step & $Q$ & $L-L^{\mathrm{ref}}$ & $Q$ & $L-L^{\mathrm{ref}}$ & $J$ & feasible \\
\midrule
3{,}000 & 27.81 & +0.016 & 42.48 & +0.038 & 35.15 & yes \\
3{,}500 & 27.99 & +0.099 & 43.19 & +0.137 & 35.59 & no \\
4{,}000 & 27.72 & +0.077 & 42.89 & +0.098 & 35.31 & yes \\
\bottomrule
\end{tabular}}
\caption{Stage-2 checkpoints on dev5 under the pre-registered rule ($Q$: text tcpBLEU-all; $L$: text LAAL of the rule's protocol, s; $J$: mean $Q$ of the two directions; feasible: $L\le L^{\mathrm{ref}}+0.1$\,s in both directions). The rule selects step 4{,}000; we use step 3{,}500 (Appendix~\ref{app:v10train}).}
\label{tab:s2sel10}
\end{table}


\paragraph{Deviation: step 3{,}500 instead of 4{,}000.} The rule stopped at step \num{v10:s2:decision} and selected step
\num{v10:s2:chosen} (Table~\ref{tab:s2sel10}). On the same dev5 numbers, and before any test job of the step-4{,}000
system had started, we chose step 3{,}500 instead, the checkpoint with the highest $J$ (\num{v10:s2:3500:J} vs.\
\num{v10:s2:4000:J}), on the principle that BLEU takes precedence over a small excess of lag. Its En$\rightarrow$Zh text
LAAL exceeds the reference by \num{v10:s2:e:3500:dL}\,s, beyond the rule's \num{v10:c:sel:tol}\,s. The decision
was made on development data only; the test sets were decoded for step 3{,}500 alone. Step 3{,}500 lies before the point
where the output starts to shorten: on dev5 (Zh$\rightarrow$En) BLEU without its brevity penalty rises over steps 3{,}000,
3{,}500 and 4{,}000 (\num{v10:c:s2:nobp3000}, \num{v10:c:s2:nobp3500}, \num{v10:c:s2:nobp4000}) while the length ratio
falls at step 4{,}000 (\num{v10:c:s2:len3500} to \num{v10:c:s2:len4000}).

\paragraph{Talkers.} The PL-EM Talker for Zh$\rightarrow$En is the one trained for the separate Zh$\rightarrow$En system
(\num{v10:c:plem:steps:z} steps); the En$\rightarrow$Zh one was trained with the same recipe, and of its checkpoints at
iterations 400, 800 and 1{,}200 the dev5 speech rule (highest ASR-BLEU among checkpoints within 0.1\,s of the last
one's speech LAAL) chose iteration \num{v10:c:plem:it:e}.

\paragraph{Other deviations.} The test decoding of stage 2 started as soon as its checkpoint was fixed, before the full
development evaluation of the joint system had finished; nothing on the test sets was used for any choice. The write bias
(Appendix~\ref{app:calib}) is the one change made after a test result had been seen.

\section{Latent Write Alignment and PL-EM: Details}
\label{app:crf}

\paragraph{Scoring.} For every window the E-step enumerates the feasible sets of Eq.~\eqref{eq:lattice} read by read; a
read with a single feasible transition (forced read) adds the same constant to every plan and is not scored. Every other
transition and order is one row: the read's input $\bm{E}^{(r)}$, the history $\kappa(A')$ and the tokens
$Y_r^o(A',A)\,\eta_r$; at the last read the wait token is masked before the softmax, as at decoding. The scores use the
training rows themselves (with the same posterior perturbation as the M-step), so the E- and M-steps see the same inputs.
Forward--backward runs in float64; no path is pruned. On the label plan ($q=1$) the M-step rows are the deadline rows
token by token: on all \num{v10:c:crf:win} windows used by a 3{,}000-step single-direction run, the rows were identical;
\num{v10:c:crf:realwin} of those windows contain a real order choice. A step scores about \num{v10:c:crf:pass1} rows
and trains on about \num{v10:c:crf:pass2} (rows sharing a prefix are merged and their coefficients added).

\paragraph{Why the factorisation is exact.} The input of read $r$ depends on the plan only through the set $A_{r-1}$
(the history is the canonical serialisation of the set, not the order in which it was written), so
$\log p_\theta$ of a plan is a sum of per-read terms that depend on $(A_{r-1},A_r,o_r)$ only, and Eq.~\eqref{eq:post} is a
chain over set states. The price is a small train--test mismatch: at decoding the history is the decoder's own order.

\paragraph{Training-time budget (plan B).} The alternative to the write bias keeps the lattice and adds the cost of
deferral to the posterior, $q_\lambda(z)\propto\exp\big(\tau\sum_r\psi_{r,o_r}-\lambda\,\mathrm{Lat}(z)\big)$, where
$\mathrm{Lat}(z)$ sums over deferred chunks their number of target units times the read interval; per window, $\lambda$ is
found by bisection so that $\mathbb{E}_{q_\lambda}[\mathrm{Lat}]=\rho\,\mathbb{E}_{q_0}[\mathrm{Lat}]$, with the single
pre-registered $\rho=0.5$. Same data, schedule and selection rule as stage 2; the rule selected step
\num{v10:pb:chosen}. Table~\ref{tab:calib10} (bottom) and Table~\ref{tab:drift10} give its results.

\begin{table}[t]
\centering\small
\setlength{\tabcolsep}{3pt}
\resizebox{\columnwidth}{!}{\begin{tabular}{lrrr}
\toprule
Training posterior (steps 1--3{,}500) & v2$'$ & $\rho=0.5$ \\
\midrule
label-path mass $q(\text{label})$, mean & 0.58 & 0.72 \\
\quad first 100 / last 500 steps & 0.66 / 0.56 & 0.77 / 0.70 \\
$\mathbb{E}_q$[deferred chunks] per window & 0.86 & 0.43 \\
\quad first 100 / last 500 steps & 0.60 / 0.91 & 0.30 / 0.46 \\
$\mathbb{E}_{q_0}$[deferral cost] per window (unit$\cdot$s) & -- & 2.89 \\
\quad first 100 / last 500 steps & -- & 2.37 / 3.37 \\
\bottomrule
\end{tabular}}
\vspace{4pt}
\resizebox{\columnwidth}{!}{\begin{tabular}{lrrrr}
\toprule
Decoding (full dev) & B & v2$'$ $\bias{=}0$ & v2$'$ $\bias{=}{-}1$ \\
\midrule
Zh$\rightarrow$En: words before unit end (\%) & 39.7 & 31.6 & 38.3 \\
\quad speech start offset (s) & 2.95 & 3.41 & 2.99 \\
\quad text LAAL (s) & 2.95 & 3.37 & 3.03 \\
En$\rightarrow$Zh: words before unit end (\%) & 55.3 & 49.1 & 53.9 \\
\quad speech start offset (s) & 2.36 & 2.74 & 2.42 \\
\quad text LAAL (s) & 2.32 & 2.70 & 2.42 \\
\bottomrule
\end{tabular}}
\caption{Late-write drift. Top: posterior statistics logged at every stage-2 training step (about 9.6 windows per step); $\rho=0.5$ is the run with an E-step latency budget (its $\mathbb{E}_{q_0}$ is the cost under the unconstrained posterior of the current model). Bottom: full development set, strict baseline (B) and the latent-write-aligned Thinker without and with the write bias.}
\label{tab:drift10}
\end{table}


\subsection{PL-EM}
\label{app:plem}

\paragraph{Candidates and fidelity.} For every block, the Talker samples 8 candidates: 4 at the rate the rule-based
Talker would use and 4 at two rate steps faster. $f(c)=-\min(\text{CTC loss per character},10)$ under an independent
wav2vec2 CTC recogniser of the target language. The forward recursion over the backlog,
$\alpha_{i+1}(h')=\log\sum_{(h,c):T_i(h,c)=h'}e^{\alpha_i(h)+\phi_i(h,c)}$, and the backward messages
$\beta_i(h)=\log\sum_c e^{\phi_i(h,c)+\beta_{i+1}(T_i(h,c))}$ agree with brute-force enumeration on small streams.

\paragraph{Weights and objective.} With $h_i^{\mathrm{rule}}$ the backlog the rule-based path reached, block $i$'s
candidates get weights $w_{i,c}\propto\exp\big(\phi_i(h_i^{\mathrm{rule}},c)+\beta_{i+1}(T_i(h_i^{\mathrm{rule}},c))\big)$
(candidates with $w<\num{v10:c:plem:minq}$ dropped). The Talker, conditioned on a backlog token, minimises the weighted
cross-entropy of the candidates' speech codes, plus \num{v10:c:plem:wrate} times the cross-entropy of its rate token
(\num{v10:c:plem:rates} rates) against the posterior's rate marginal, plus \num{v10:c:plem:wkl} times the KL divergence
from the frozen T1c Talker on the semantic codes and the end token. All Talker parameters are trained.

\paragraph{Data.} Zh$\rightarrow$En: \num{v10:c:plem:meet:z} AMI training meetings decoded by the separately trained
Zh$\rightarrow$En latent-write Thinker, \num{v10:c:plem:blocks:z} blocks, \num{v10:c:plem:rows:z} weighted candidate rows;
budget: the rule path's own delay minus \num{v10:c:plem:budgetoff}\,s, $\lambda^\star=$ \num{v10:c:plem:lam:z}.
En$\rightarrow$Zh: the same recipe, \num{v10:c:plem:rows:e} rows, $\lambda^\star=$ \num{v10:c:plem:lam:e}.

\begin{table}[t]
\centering\small
\setlength{\tabcolsep}{3pt}
\resizebox{\columnwidth}{!}{\begin{tabular}{lrrrrrr}
\toprule
 & \multicolumn{3}{c}{AMI Zh$\rightarrow$En} & \multicolumn{3}{c}{AliMeeting} \\
\cmidrule(lr){2-4}\cmidrule(lr){5-7}
Talker & ASR-BLEU & LAAL$_\text{spr.}$ & EndOff. & ASR-BLEU & LAAL$_\text{spr.}$ & EndOff. \\
\midrule
rule (queue rate, 2-best CTC) & 20.52 & 4.00 & 3.65 & 15.83 & 4.73 & 4.42 \\
PA Talker & 20.64 & 3.92 & 3.42 & 16.26 & 4.61 & 4.08 \\
rl + PA Talker & 20.68 & 3.90 & 3.36 & 16.22 & 4.57 & 4.00 \\
\bottomrule
\end{tabular}}
\caption{Speech layer alone, on a separately trained single-direction Zh$\rightarrow$En write-aligned Thinker (same text in all rows: tcpBLEU 25.37 / 18.14); test sets, two TTS seeds, speech LAAL and EndOffset in seconds (chunk protocol). rl: read-level merging. The rule-based layer here already uses two samples per block with CTC selection and a queue-dependent rate.}
\label{tab:talker10}
\end{table}


\section{Write-Bias Calibration and the Run Without It}
\label{app:calib}

\paragraph{When and how.} The first test decoding of the latent-write-aligned system used no bias. It showed that the
system started speaking later than the strict baseline, most on AliMeeting (Table~\ref{tab:delta0}). We then added
Eq.~\eqref{eq:bias} and calibrated it on dev5 only: candidates $\{-1,-2,-3,-5\}$ in both directions, rule fixed before
decoding (text LAAL closest to the strict baseline's), no finer grid. Both directions gave $\bias=-1$
(Table~\ref{tab:calib10}): text LAAL \num{v10:cal:z:m1-b:laal} / \num{v10:cal:e:m1-b:laal}\,s from the baseline's,
tcpBLEU \num{v10:cal:z:m1-b:bleu} / \num{v10:cal:e:m1-b:bleu}. Larger biases cut the translation short
(\num{v10:cal:z:m3:bleu} tcpBLEU at $-3$ in Zh$\rightarrow$En). The value was confirmed once on the full development set
(Table~\ref{tab:abl10}), and each test set was decoded once with it. The decision to introduce a bias at all was
influenced by the first test run; its value was not.

\begin{table}[t]
\centering\small
\setlength{\tabcolsep}{4pt}
\resizebox{\columnwidth}{!}{\begin{tabular}{lrrrr}
\toprule
 & \multicolumn{2}{c}{Zh$\rightarrow$En} & \multicolumn{2}{c}{En$\rightarrow$Zh} \\
\cmidrule(lr){2-3}\cmidrule(lr){4-5}
dev5, text only & tcpBLEU & LAAL & tcpBLEU & LAAL \\
\midrule
deadline baseline (B) & 28.11 & 3.07 & 41.44 & 2.35 \\
write alignment, $\bias=0$ & 27.99 & 3.54 & 43.19 & 2.81 \\
\quad $\bias=-1$ (chosen) & 28.25 & 3.16 & 42.29 & 2.45 \\
\quad $\bias=-2$ & 25.25 & 2.55 & 39.16 & 2.17 \\
\quad $\bias=-3$ & 14.66 & 1.95 & 33.93 & 1.90 \\
\quad $\bias=-5$ & 6.51 & 1.30 & 17.28 & 1.39 \\
\midrule
E-step budget $\xi=0.5$, step 3{,}500 (rule) & 27.92 & 3.45 & 42.90 & 2.68 \\
\quad step 5{,}000 & 28.20 & 3.37 & 42.99 & 2.59 \\
\bottomrule
\end{tabular}}
\caption{Write-bias calibration on the five development meetings (dev5; text decoding, chunk-protocol text LAAL in seconds). The rule, fixed before decoding: per direction, the $\bias$ whose text LAAL is closest to the deadline baseline's; both directions gave $\bias=-1$, which was then checked on the full development set (which contains dev5, so the check is not independent). Bottom: write alignment retrained with an E-step latency budget ($\xi=0.5$, the single value fixed in advance), checkpoint chosen by the same rule as stage 2, $\bias=0$.}
\label{tab:calib10}
\end{table}

\begin{table}[t]
\centering\small
\setlength{\tabcolsep}{2.5pt}
\resizebox{\columnwidth}{!}{\begin{tabular}{llrrrrr}
\toprule
Set & System & tcpBLEU & ov & ASR-BLEU & LAAL$_\text{spr.}$ & EndOff. \\
\midrule
AMI Zh$\rightarrow$En & B & 25.02 & 19.97 & 18.07 & 4.01 & 4.08 \\
 & $\bias=0$ (first run) & 25.09 & 19.96 & 20.48 & 3.99 & 3.42 \\
 & $\bias=-1$ & 25.34 & 20.53 & 20.45 & 3.69 & 3.20 \\
\midrule
AliMeeting & B & 21.43 & 16.91 & 16.73 & 4.28 & 4.57 \\
 & $\bias=0$ (first run) & 18.99 & 15.63 & 17.37 & 4.85 & 4.23 \\
 & $\bias=-1$ & 21.07 & 17.38 & 17.84 & 3.91 & 3.45 \\
\midrule
AMI En$\rightarrow$Zh & B & 41.15 & 36.16 & 31.65 & 3.21 & 3.48 \\
 & $\bias=0$ (first run) & 41.85 & 36.18 & 32.41 & 3.34 & 3.21 \\
 & $\bias=-1$ & 41.74 & 36.48 & 32.34 & 3.05 & 2.99 \\
\bottomrule
\end{tabular}}
\caption{Test results of the latent-write-aligned system without the write bias ($\bias=0$), the run that preceded the introduction of $\bias$, next to the deadline baseline (B) and the final system. Same speech layer (rl + PA Talker) for both $\bias$ values; chunk protocol.}
\label{tab:delta0}
\end{table}


\section{Latency Protocol and Delay Decomposition}
\label{app:chunk}

\paragraph{Chunk protocol.} The low-latency diarizer emits frames in blocks of nine; frame $f$ (80\,ms) is known at
$T(f)=(9\lfloor f/9\rfloor+13)\cdot 80\,\text{ms}\le f\cdot80\,\text{ms}+1.04$\,s. A window's opening is released at
$T(\text{first active frame})$ and its end at $T(\text{last frame of the closing gap})$; commits before the opening are
delayed to it, and those of the last read to the release of the end. The v9 protocol released both at frame time plus
1.04\,s. This is a bookkeeping correction applied to every system alike: re-scoring \num{v10:c:chunk:nsys} earlier
systems under both protocols left their ranking unchanged and lowered speech LAAL by
\num{v10:c:chunk:shiftlo}--\num{v10:c:chunk:shifthi}\,s (openings are released \num{v10:c:chunk:open}\,ms and ends
\num{v10:c:chunk:close}\,ms earlier on average). Systems are compared only under the same protocol.

\paragraph{Speech LAAL.} For a reference unit with source duration $|X|$ and $|Y|$ heard target words (words or Chinese
characters) with delays $d_1\le\dots\le d_{|Y|}$ from the unit's start, $\mathrm{LAAL}=\frac{1}{\tau'}\sum_{k=1}^{\tau'}
\big(d_k-(k-1)\frac{|X|}{\max(|Y|,|Y^\ast|)}\big)$, with $|Y^\ast|$ the reference length and $\tau'$ the first $k$ with
$d_k\ge|X|$~\citep{papi2022over}. A word's delay is the moment it is heard: block $i$ of a stream starts playing at
$\nu_i+h_i$ and its $W_i$ words are spread evenly over its duration. Units, matching and collars are those of tcpBLEU;
each speaker has its own stream. Table~\ref{tab:m2laal10} also gives the convention of single-speaker work in which every
word of a block is charged at the block's start, which cannot see how long a block takes to say.

\begin{table}[t]
\centering\small
\setlength{\tabcolsep}{4pt}
\begin{tabular}{lrrrr}
\toprule
 & \multicolumn{2}{c}{block start} & \multicolumn{2}{c}{word heard (ours)} \\
\cmidrule(lr){2-3}\cmidrule(lr){4-5}
Set & B & \sys{} & B & \sys{} \\
\midrule
dev Zh$\rightarrow$En & 3.10 & 3.06 & 3.97 & 3.61 \\
dev En$\rightarrow$Zh & 2.49 & 2.48 & 3.12 & 2.96 \\
test AMI Zh$\rightarrow$En & 3.12 & 3.39 & 4.01 & 3.99 \\
test AliMeeting & 3.28 & 4.01 & 4.28 & 4.85 \\
test AMI En$\rightarrow$Zh & 2.54 & 2.80 & 3.21 & 3.34 \\
\bottomrule
\end{tabular}
\caption{Speech LAAL with the single-speaker convention (every word of a block is delayed until the block \emph{starts} playing) next to ours (each word at the moment it is heard). B: strict baseline; \sys{}: $\bias=-1$ on dev, $\bias=0$ on test (the $\bias=-1$ test runs were not rescored with this convention). Chunk protocol, mean of two TTS seeds, seconds.}
\label{tab:m2laal10}
\end{table}


\subsection{Where the delay comes from}
\label{app:decomp}

On the separate Zh$\rightarrow$En system, speech LAAL (\num{v10:c:dec:speech}\,s) is the raw commit delay
(\num{v10:c:dec:raw}\,s) plus \num{v10:c:dec:det}\,s of detector wait and \num{v10:c:dec:play}\,s of playback, and text
LAAL after the detector (\num{v10:c:dec:online}\,s) equals the mean reference-unit duration; at the utterance level the
system interprets consecutively.

Table~\ref{tab:decomp10} splits EndOffset into its parts. The speech layer of the baseline (fixed rate, every committed
block in full) produces more speech than the source holds, so each stream accumulates a queue; PL-EM with read-level
merging removes most of it. On the separate Zh$\rightarrow$En system (development set,
\num{v10:c:dec:units} units, mean duration \num{v10:c:dec:unitdur}\,s, median \num{v10:c:dec:unitp50}\,s,
\num{v10:c:dec:short}\% shorter than 2\,s) we replayed the same text with every segment committed at the
first read after its reference unit ends plus a lag: with 1\,s reads and 0.3\,s lag speech LAAL falls by
\num{v10:c:ub:l03}\,s, with 0.5\,s reads by \num{v10:c:ub:h03}\,s (0.3\,s lag) and \num{v10:c:ub:h0}\,s (no lag). That is
the most a better write decision at utterance level can give; the rest of the gap to single-speaker systems needs
writing within the utterance.

\begin{table}[t]
\centering\small
\setlength{\tabcolsep}{2.5pt}
\resizebox{\columnwidth}{!}{\begin{tabular}{llrrrrrrr}
\toprule
Set & Sys. & start & model & detector & queue & last block & EndOff. & load \\
\midrule
dev Zh$\rightarrow$En & B & 2.95 & 0.74 & 0.61 & 0.64 & 2.16 & 4.14 & 1.41 \\
 & B+PL & 2.90 & 0.75 & 0.61 & 0.12 & 1.62 & 3.10 & 0.94 \\
 & \sys{} & 2.99 & 0.79 & 0.61 & 0.13 & 1.65 & 3.18 & 0.95 \\
\midrule
dev En$\rightarrow$Zh & B & 2.36 & 0.76 & 0.55 & 0.59 & 1.54 & 3.45 & 1.33 \\
 & B+PL & 2.32 & 0.78 & 0.55 & 0.14 & 1.37 & 2.83 & 1.00 \\
 & \sys{} & 2.42 & 0.83 & 0.56 & 0.14 & 1.40 & 2.93 & 1.00 \\
\midrule
AMI Zh$\rightarrow$En & B & 2.99 & 0.67 & 0.66 & 0.65 & 2.10 & 4.08 & 1.37 \\
 & B+PL & 2.96 & 0.68 & 0.66 & 0.14 & 1.64 & 3.13 & 0.94 \\
 & \sys{} & 3.04 & 0.71 & 0.67 & 0.14 & 1.67 & 3.20 & 0.95 \\
\midrule
AliMeeting & B & 3.05 & 0.88 & 0.40 & 1.33 & 1.96 & 4.57 & 1.56 \\
 & B+PL & 2.92 & 0.95 & 0.40 & 0.33 & 1.55 & 3.23 & 1.05 \\
 & \sys{} & 3.20 & 1.02 & 0.40 & 0.35 & 1.67 & 3.45 & 1.06 \\
\midrule
AMI En$\rightarrow$Zh & B & 2.43 & 0.69 & 0.64 & 0.54 & 1.61 & 3.48 & 1.28 \\
 & B+PL & 2.39 & 0.71 & 0.64 & 0.14 & 1.42 & 2.91 & 0.99 \\
 & \sys{} & 2.50 & 0.73 & 0.65 & 0.14 & 1.46 & 2.99 & 0.99 \\
\bottomrule
\end{tabular}}
\caption{Where the speech delay of a source unit comes from (chunk protocol, mean over units and two TTS seeds, seconds). EndOffset = model lag (last commit after the unit's end) + detector wait (online window release) + queue wait (earlier speech of the same stream) + duration of the last block. start: speech start offset from the unit's start. load: synthesised speech per second of source speech. B+PL: strict baseline text with rl + PL-EM.}
\label{tab:decomp10}
\end{table}


\section{Further v10 Results}

\subsection{AliMeeting: brevity penalty and duplicates}
\label{app:dup}

The overlap scorer pools the segments of all online windows and attaches each to the reference unit it overlaps, so a
unit translated again in a neighbouring window keeps both translations. On the AliMeeting test set the baseline has
\num{v10:c:dup:ali:b:words} such duplicated words (\num{v10:c:dup:ali:b:share}\% of its \num{v10:c:dup:ali:b:att}
attached words), the $\bias=0$ system \num{v10:c:dup:ali:m0:words}; without them the system writes
\num{v10:c:dup:ali:gain}\% more non-duplicated words (\num{v10:c:dup:ali:m0:nondup} vs.\ \num{v10:c:dup:ali:b:nondup}). Keeping
only the longest translation of each duplicated unit, tcpBLEU is \num{v10:c:dd:ali:b} vs.\ \num{v10:c:dd:ali:m0}
(AliMeeting), \num{v10:c:dd:ami:b} vs.\ \num{v10:c:dd:ami:m0} (AMI Zh$\rightarrow$En) and \num{v10:c:dd:devz:b} vs.\
\num{v10:c:dd:devz:m0} (development Zh$\rightarrow$En), baseline vs.\ $\bias=0$. The same drop on AliMeeting appears with
the separately trained Zh$\rightarrow$En latent-write Thinker (\num{v10:c:ali:sepA} vs.\ \num{v10:test:ali:b:all}), so it
is a property of the write plan, not of the joint model. These are diagnostics; the official tcpBLEU is the reported
metric. Table~\ref{tab:dup10} recomputes the brevity-penalty split from the score files for all three systems.

\begin{table}[t]
\centering\small
\setlength{\tabcolsep}{3pt}
\resizebox{\columnwidth}{!}{\begin{tabular}{llrrrrr}
\toprule
Set & System & BLEU & len.\ ratio & BP & BLEU/BP & dup.\ units (\%) \\
\midrule
AliMeeting & B & 21.43 & 1.001 & 1.000 & 21.43 & 18.1 \\
 & $\bias=0$ & 18.99 & 0.913 & 0.909 & 20.90 & 11.6 \\
 & $\bias=-1$ & 21.07 & 0.984 & 0.984 & 21.42 & 19.2 \\
\midrule
AMI Zh$\rightarrow$En & B & 25.02 & 0.946 & 0.944 & 26.50 & 6.5 \\
 & $\bias=0$ & 25.09 & 0.917 & 0.914 & 27.45 & 6.2 \\
 & $\bias=-1$ & 25.34 & 0.943 & 0.942 & 26.90 & 7.7 \\
\midrule
AMI En$\rightarrow$Zh & B & 41.15 & 1.005 & 1.000 & 41.15 & 7.8 \\
 & $\bias=0$ & 41.85 & 0.969 & 0.968 & 43.22 & 7.0 \\
 & $\bias=-1$ & 41.74 & 1.000 & 1.000 & 41.74 & 8.6 \\
\midrule
dev Zh$\rightarrow$En & B & 27.93 & 0.967 & 0.966 & 28.90 & 5.9 \\
 & $\bias=0$ & 27.75 & 0.931 & 0.928 & 29.90 & 5.2 \\
 & $\bias=-1$ & 28.31 & 0.968 & 0.967 & 29.28 & 6.5 \\
\bottomrule
\end{tabular}}
\caption{Text tcpBLEU split into its brevity penalty (BP) and the precision part (BLEU/BP), with the share of reference units that received two near-identical translations (dup.). Computed from the official score files; diagnostic only.}
\label{tab:dup10}
\end{table}


\subsection{Significance}
\label{app:sig10}

\begin{table}[t]
\centering\small
\setlength{\tabcolsep}{2.5pt}
\resizebox{\columnwidth}{!}{\begin{tabular}{llrrrr}
\toprule
Set & Comparison & $\Delta$ & 95\% interval & $p$ & $p_\text{Holm}$ \\
\midrule
AMI Zh$\rightarrow$En & text: \sys{} $-$ B & +0.32 & [0.01, 0.64] & 0.081 & 0.081 \\
 & speech: \sys{} $-$ B & \textbf{??} & \textbf{??} & \textbf{??} & \textbf{??} \\
 & speech: B+PL $-$ B & \textbf{??} & \textbf{??} & \textbf{??} & \textbf{??} \\
 & speech: \sys{} $-$ B+PL & \textbf{??} & \textbf{??} & \textbf{??} & \textbf{??} \\
\midrule
AliMeeting & text: \sys{} $-$ B & \textbf{??} & \textbf{??} & \textbf{??} & \textbf{??} \\
 & speech: \sys{} $-$ B & \textbf{??} & \textbf{??} & \textbf{??} & \textbf{??} \\
 & speech: B+PL $-$ B & \textbf{??} & \textbf{??} & \textbf{??} & \textbf{??} \\
 & speech: \sys{} $-$ B+PL & \textbf{??} & \textbf{??} & \textbf{??} & \textbf{??} \\
\midrule
AMI En$\rightarrow$Zh & text: \sys{} $-$ B & +0.58 & [0.12, 0.90] & 0.013 & 0.027 \\
 & speech: \sys{} $-$ B & \textbf{??} & \textbf{??} & \textbf{??} & \textbf{??} \\
 & speech: B+PL $-$ B & \textbf{??} & \textbf{??} & \textbf{??} & \textbf{??} \\
 & speech: \sys{} $-$ B+PL & \textbf{??} & \textbf{??} & \textbf{??} & \textbf{??} \\
\bottomrule
\end{tabular}}
\caption{Paired tests over meetings on the test sets (16 AMI, 20 AliMeeting meetings): exact two-sided sign-flip permutation $p$ over all $2^M$ flips, 95\% meeting-bootstrap interval (10{,}000 resamples), Holm over the 2 tests. Speech pools the two TTS seeds per meeting. B: strict baseline; B+PL: its text with rl + PL-EM; \sys{}: v10 with $\bias=-1$. Text of B+PL equals B by construction.}
\label{tab:sig10}
\end{table}


Table~\ref{tab:sig10} lists the paired tests on the test sets. On the development set, \sys{} v10 against the strict
baseline changes text tcpBLEU by \num{v10:sig:dev:z:text:m-b:diff} (\num{v10:sig:dev:z:text:m-b:ci}) /
\num{v10:sig:dev:e:text:m-b:diff} (\num{v10:sig:dev:e:text:m-b:ci}) and ASR-BLEU by
\num{v10:sig:dev:z:speech:m-b:diff} (\num{v10:sig:dev:z:speech:m-b:ci}) / \num{v10:sig:dev:e:speech:m-b:diff}
(\num{v10:sig:dev:e:speech:m-b:ci}), Zh$\rightarrow$En / En$\rightarrow$Zh.

\subsection{External system: SeamlessStreaming}
\label{app:e7}

\ifEseven
SeamlessStreaming~\citep{seamless2023} (official agent and default settings, inference only) reads each of our online
test windows; each output segment is attributed to a diarizer slot by the causal relabelling rule of
\S\ref{sec:eval-overlap}, and text, speech and latency are scored with our tools under the chunk protocol. It thus gets our
segmentation and speaker stream, which it does not have on its own. Results: Table~\ref{tab:main10}.
\else
% E7 slot: filled when reports/mist_w3po_v1/e7_seamless_multispk_results.zh-CN.md and
% eval_outputs/mist_w3po_v1/decode/jtest_<set>_seamless_default/score_text.json exist (rerun make_tables_v10.py).
SeamlessStreaming~\citep{seamless2023} on our online test windows (official agent and defaults, inference only, causal
diarizer relabelling of its segments) was registered as an external baseline; its run had not finished for this version.
\fi

\section{The v9 System Under the v9 Protocol}
\label{app:v9proto}

The results of this section concern the v9 deadline-trained \sys{}-Simul (one model per direction), the v9 online
windows (1\,s closing gap), the v9 latency protocol (release at frame time plus 1.04\,s; reads queued on one decoder
stream at the measured single-stream read time) and the listener-side lag of the v8 queue model. They are not comparable
with the v10 numbers.

\begin{table*}[t]
\centering\small
\setlength{\tabcolsep}{4pt}
\resizebox{\textwidth}{!}{\begin{tabular}{lccccc@{\hskip 10pt}ccccc}
\toprule
 & \multicolumn{5}{c}{En$\rightarrow$Zh} & \multicolumn{5}{c}{Zh$\rightarrow$En} \\
\cmidrule(lr){2-6}\cmidrule(lr){7-11}
System & LAAL\,s & cpBLEU$_\text{sess}$ & BLEU$_\text{blind}$ & Spk & DER & LAAL\,s & cpBLEU$_\text{sess}$ & BLEU$_\text{blind}$ & Spk & DER \\
\midrule
\multicolumn{11}{l}{\emph{Offline}} \\
Cascade, same-data translator & -- & 36.29 & 37.01 & 88.2 & 25.8 & -- & 22.36 & 24.85 & 83.6 & 15.5 \\
Cascade, turn-trained translator & -- & 35.27 & 35.27 & 79.3 & 37.8 & -- & 23.69 & 25.32 & 86.4 & 16.9 \\
Cascade, turns + 2\,s context & -- & 23.75 & 25.02 & 88.3 & 25.9 & -- & 16.88 & 19.53 & 83.5 & 15.6 \\
\sys{} w/o speaker stream & -- & 21.17 & 43.40 & 84.8 & 54.5 & -- & 12.20 & 28.13 & 84.9 & 57.4 \\
\quad + diarizer relabelling & -- & 42.06 & 43.40 & 82.2 & 13.9 & -- & 24.66 & 28.13 & 82.0 & 15.7 \\
\sys{} (offline) & -- & 42.81$^\dagger$ & 43.33 & 89.5 & 14.0 & -- & 25.86$^\dagger$ & 28.43 & 91.3 & 16.0 \\
\addlinespace
\multicolumn{11}{l}{\emph{Simultaneous ($\Delta{=}1$\,s)}} \\
\sys{} (offline) + LA2 & 2.37 & 31.64 & 32.74 & 87.8 & 15.2 & 2.50 & 24.81 & 27.67 & 90.7 & 17.3 \\
\sys{}-Simul model + LA2 & 3.07 & 39.21 & 39.64 & 88.4 & 12.0 & 2.97 & 24.86 & 27.21 & 91.7 & 13.2 \\
Streaming cascade (turns + LA2) & 2.50 & 35.16 & 35.84 & 87.4 & 26.1 & 2.70 & 21.55 & 23.68 & 80.3 & 16.2 \\
\sys{}-Simul & 2.40 & 43.08$\pm$0.45$^{3}$ & 43.24 & 90.4 & 11.3 & 3.04 & 23.55$\pm$1.07$^{2}$ & 25.88 & 93.1 & 12.3 \\
\addlinespace
\multicolumn{11}{l}{\emph{Oracle input (reference turns or activity), same translator}} \\
Cascade, reference turns (oracle) & -- & 37.43 & 38.32 & 99.7 & 0.5 & -- & 24.95 & 26.28 & 100.0 & 1.0 \\
\sys{} (offline), oracle stream & -- & 42.32 & 42.95 & 92.2 & 14.2 & -- & 27.32 & 28.37 & 95.9 & 11.6 \\
\bottomrule
\end{tabular}}
\caption{Speaker-attributed translation (text) on the reference windows of the AMI test meetings (2{,}375 En$\rightarrow$Zh and 3{,}317 Zh$\rightarrow$En windows, 16 meetings). One bidirectional \sys{} (offline) serves both directions; \sys{}-Simul is trained per direction and reported as the mean $\pm$ s.d.\ over its training seeds (superscript: number of seeds; per seed in Table~\ref{tab:seeds7}). cpBLEU$_\text{sess}$ (primary) matches output speakers to reference speakers once per meeting; BLEU$_\text{blind}$ ignores speakers. Spk: windows with the right number of speakers (\%); DER (\%) of the output segments. LAAL includes the diarizer's 1.04\,s look-ahead and is computation-unaware. $^\dagger$: better than the strongest non-oracle baseline of its block (exact permutation test over meetings $p<0.05$ and the decoder-noise-widened interval excludes 0; Table~\ref{tab:sig7}).}
\label{tab:main7full}
\end{table*}


\paragraph{Speaker-blind Simul control.} Online, a speaker-blind Simul model trained identically and relabelled causally
(pre-registered E-ov1/E-ov2, Table~\ref{tab:reg8}) trails \sys{}-Simul on overlapped units by
\num{v8:reg:ov1-en2zh:diff} (En$\rightarrow$Zh) and \num{v8:reg:ov1-zh2en:diff} (Zh$\rightarrow$En) tcpBLEU, but ties on
AliMeeting (\num{v8:reg:ov1-ali:diff}), where it is even ahead overall; concurrent-speaker recall is higher for
\sys{}-Simul in all three sets (by \num{v8:reg:ov2-en2zh:diff} / \num{v8:reg:ov2-zh2en:diff} /
\num{v8:reg:ov2-ali:diff}). Adding the cascade's translation of short turns \sys{} left out raises concurrent recall to
\num{v8:ov:en2zh:C2_O:cr}\% and tcpBLEU to \num{v8:ov:en2zh:C2_O:all} (Appendix~\ref{app:v8}).

\paragraph{Commit policy and latency.}
\label{sec:res-simul}
\label{sec:res-latency}

\begin{table}[t]
\centering\small
\setlength{\tabcolsep}{2.2pt}
\resizebox{\columnwidth}{!}{\begin{tabular}{llrrrrrrr}
\toprule
 & & & \multicolumn{2}{c}{text lag p50 (s)} & \multicolumn{3}{c}{per-speaker channel (s)} & 1 ch. \\
\cmidrule(lr){4-5}\cmidrule(lr){6-8}
Set & System & RTF & overlap & single & ovl. p50 & single p50 & p90 & p50 \\
\midrule
En$\rightarrow$Zh & \sys{}-Simul & 0.86 & 1.59 & 1.96 & 6.8 & 7.7 & 12.2 & 142.4 \\
 & \sys{} + LA2 & 1.73 & 3.01 & 2.02 & 874.1 & 821.4 & 1688.6 & 903.0 \\
Zh$\rightarrow$En & \sys{}-Simul & 0.72 & 1.73 & 1.95 & 7.5 & 8.5 & 15.1 & 26.4 \\
 & \sys{} + LA2 & 1.45 & 2.99 & 1.98 & 613.7 & 445.9 & 1009.6 & 500.1 \\
AliMeeting & \sys{}-Simul & 0.98 & 1.47 & 1.82 & 33.7 & 20.2 & 314.1 & 604.6 \\
 & \sys{} + LA2 & 1.98 & 4.09 & 1.99 & 1045.8 & 990.2 & 1783.9 & 1252.5 \\
\bottomrule
\end{tabular}}
\caption{Whole meetings with queued reads. RTF: read computation per second of meeting ($>1$: falls behind). Text
lag: commit of a unit's last segment minus the unit's end. Per-speaker channel and one channel (1 ch.): completion lag
of the speech heard.}
\label{tab:lat8}
\end{table}


On reference windows LocalAgreement reaches \num{v7:tx:en2zh:O2_LA2_d1:cp} cpBLEU$_\text{sess}$ (En$\rightarrow$Zh) on
the offline model and \num{v7:tx:en2zh:S2dl_LA2_full:cp} on the prefix-trained one; \sys{}-Simul's wait policy
(\num{v7:tx:en2zh:S2dl_d1:cp}) is ahead of the latter within noise and 0.75\,s earlier. On overlapped units the same comparison favours the learned policy
(\num{v8:ov:en2zh:S2dl_d1:ov} vs.\ \num{v8:ov:en2zh:S2dl_LA2_full:ov}; Table~\ref{tab:overlap8}): LocalAgreement waits for
two consecutive hypotheses to agree, which overlapping speakers rarely let happen. Its real cost is computation:
LocalAgreement decodes to the end at every read (\num{v6:lat:en2zh:O2_LA2:read_ms}\,ms per read, against
\num{v6:lat:en2zh:S2dl:read_ms}\,ms). With the reads of a meeting queued on one decoder stream (Table~\ref{tab:lat8}) its
real-time factor is \num{v8:lat:en2zh:O2_LA2:single:rtf}--\num{v8:lat:ali:O2_LA2:single:rtf}, so it falls further behind
the longer the meeting runs, whereas \sys{}-Simul's is \num{v8:lat:en2zh:S2dl:single:rtf} (En$\rightarrow$Zh),
\num{v8:lat:zh2en:S2dl:single:rtf} (Zh$\rightarrow$En) and \num{v8:lat:ali:S2dl:single:rtf} (AliMeeting) and its median
per-speaker listener lag rises only from \num{v8:lat:en2zh:S2dl:ps_noq:p50} to \num{v8:lat:en2zh:S2dl:ps:p50}\,s in
En$\rightarrow$Zh when the reads are queued. Overlap does not slow it down: its text lag on overlapped units is no larger
than on single-speaker ones (\num{v8:ov:en2zh:on_S2dl:lagov} vs.\ \num{v8:ov:en2zh:on_S2dl:lag1spk}\,s; E-lat holds in
all three sets, Table~\ref{tab:reg8}), and \num{v8:ov:en2zh:on_S2dl:pre}\% of its output words are committed before the
utterance they translate has ended. Across training seeds \sys{}-Simul's whole-meeting tcpBLEU is \num{v8:seed:en2zh:all}
(En$\rightarrow$Zh, \num{v8:seed:en2zh:n} seeds); on the meeting subsets a fixed 2\,s lag reaches the same or higher quality
(behind on AMI overlapped units, ahead on AliMeeting) at about 1\,s more lag on overlapped units (Appendix~\ref{app:v8}).

\paragraph{Whole meetings, run online.}
\label{sec:res-online}

We run the simultaneous systems over all 16 AMI and 20 AliMeeting test meetings with online windows
(\S\ref{sec:method-online}; Table~\ref{tab:online7}): no reference boundary, all audio processed, commits delayed until a
window's opening and end are known; the speaker stream is causal (Appendix~\ref{app:causal}). Against
LocalAgreement on the offline model (pre-registered E-on1--3) \sys{}-Simul is ahead by \num{v7:reg:E-on1:diff} in
En$\rightarrow$Zh ($p_\text{Holm}$ \num{v7:reg:E-on1:holm}) and ties in Zh$\rightarrow$En (\num{v7:reg:E-on2:diff}) and on
AliMeeting (\num{v7:reg:E-on3:diff}). Scored on all En$\rightarrow$Zh units instead of the
\num{v6:filt:en2zh:all:dur}\% the reference windows keep, the lead grows
(\num{v8:full:en2zh:on_S2dl:all} vs.\ \num{v8:full:en2zh:on_O2_LA2:all}; overlapped units \num{v8:full:en2zh:on_S2dl:ov}
vs.\ \num{v8:full:en2zh:on_O2_LA2:ov}), so the filtered references did not flatter \sys{}. Against LocalAgreement on
the same prefix-trained model (E-pol) its median text lag is \num{v8:reg:pol-en2zh:diff} / \num{v8:reg:pol-zh2en:diff} /
\num{v8:reg:pol-ali:diff}\,s shorter per meeting (Table~\ref{tab:reg8}) at tied quality (Table~\ref{tab:ovmain8}).

\paragraph{What a listener hears.} One interpreter channel, which serialises overlapping speakers, cannot keep up: the
median block of \sys{}-Simul is heard \num{v8:lat:en2zh:S2dl:single:p50}\,s after its source ended in
En$\rightarrow$Zh and \num{v8:lat:ali:S2dl:single:p50}\,s on AliMeeting. Captions for short overlapped interjections and skipping blocks
more than 10\,s late bound it (\num{v8:b7:en2zh:d:p50}\,s) only by speaking just \num{v8:b7:en2zh:d:cov}\% of the units (the rest become captions)
(Table~\ref{tab:sched8}). One channel per speaker keeps overlap as in the meeting: median \num{v8:lat:en2zh:S2dl:ps:p50}\,s
(90th percentile \num{v8:lat:en2zh:S2dl:ps:p90}\,s), Zh$\rightarrow$En \num{v8:lat:zh2en:S2dl:ps:p50}\,s, AliMeeting
\num{v8:lat:ali:S2dl:ps:p50}\,s.


\section{AliMeeting Human References (v9)}
\label{app:v9}

The DiariST release~\citep{yang2024diarist} provides human English translations of the AliMeeting test set. Its
\num{v9:ali:align:rows} rows all match our test utterances by meeting, speaker, times and Chinese text
(\num{v9:ali:align:matched} matched), which gives human references for every one of the \num{v9:ali:align:scored} scored
utterances. With the machine-translated references used before v9, \sys{} (offline) against the same-data cascade is
\num{v9:x:ali:O2-C-ft2:std:all:diff} tcpBLEU overall (interval \num{v9:x:ali:O2-C-ft2:std:all:ci}) and
\num{v9:x:ali:O2-C-ft2:std:ov:diff} on overlapped units (\num{v9:x:ali:O2-C-ft2:std:ov:ci}); with the human references
the corresponding numbers are those of \S\ref{sec:res-real}. Table~\ref{tab:reg9} lists the pre-registered v9 tests.

\begin{table*}[t]
\centering\small
\setlength{\tabcolsep}{3pt}
\resizebox{\textwidth}{!}{\begin{tabular}{lllrrrl}
\toprule
Test & Set & Comparison (whole meetings, online) & $\Delta$ & 95\% interval & $p_\text{Holm}$ & \\
\midrule
E-pareto-en2zh & En$\rightarrow$Zh & \sys{}-Simul not dominated by LA2 / lag 1, 2, 3\,s / block 2\,s (quality, text lag) & -- & -- & -- & not run \\
E-ovcap1-en2zh & En$\rightarrow$Zh & tcpBLEU-ov, \sys{}-Simul $>$ D1 and D2 & -- & -- & -- & not run \\
E-ovcap2-en2zh & En$\rightarrow$Zh & concurrent recall, \sys{}-Simul $>$ D1 and D2 & -- & -- & -- & not run \\
E-pareto-zh2en & Zh$\rightarrow$En & \sys{}-Simul not dominated by LA2 / lag 1, 2, 3\,s / block 2\,s (quality, text lag) & -- & -- & -- & not run \\
E-ovcap1-zh2en & Zh$\rightarrow$En & tcpBLEU-ov, \sys{}-Simul $>$ D1 and D2 & -- & -- & -- & not run \\
E-ovcap2-zh2en & Zh$\rightarrow$En & concurrent recall, \sys{}-Simul $>$ D1 and D2 & -- & -- & -- & not run \\
E-pareto-ali & AliMeeting & \sys{}-Simul not dominated by LA2 / lag 1, 2, 3\,s / block 2\,s (quality, text lag) & -- & -- & -- & not run \\
E-ovcap1-ali & AliMeeting & tcpBLEU-ov, \sys{}-Simul $>$ D1 and D2 & -- & -- & -- & not run \\
E-ovcap2-ali & AliMeeting & concurrent recall, \sys{}-Simul $>$ D1 and D2 & -- & -- & -- & not run \\
E-far & AMI SDM & SDM far field: tcpBLEU-ov, \sys{}-Simul $>$ causal relabel & -- & -- & -- & not run \\
E-ali-hum & AliMeeting & AliMeeting, human references: tcpBLEU-ov, \sys{}-Simul $>$ causal relabel & 0.61 & [-3.74, 4.18] & 0.387 & tie / fails \\
\bottomrule
\end{tabular}}
\caption{The eleven confirmatory tests of prereg\_v9 (committed before any v9 decode), Holm-corrected together. A test
holds when $p_\text{Holm}<0.05$ and, for differences, the interval excludes 0. E-pareto: for every competitor, \sys{}-Simul
must be better on quality or on text lag (two-dimension Bonferroni); the set's $p$ is the largest over competitors.
LA2 and lag 2\,s: all meetings; lag 1 and 3\,s, block 2\,s and D2: the v7 meeting subset (8 AMI / 10 AliMeeting
meetings; prereg\_v9 amendment 2), with \sys{}-Simul scored on the same meetings.
E-ovcap: the larger $p$ of the comparisons with D1 and D2; $\Delta$ and the interval are the less favourable of the two.}
\label{tab:reg9}
\end{table*}

